\documentclass[11pt,a4paper,landscape]{article}
\usepackage[T1]{fontenc}
\usepackage[utf8]{inputenc}
\usepackage[provide=*,lithuanian]{babel}
\usepackage{lmodern}
\usepackage{amsmath,amssymb}
\usepackage[margin=0pt]{geometry}
\usepackage{tikz}
\pagestyle{empty}
\setlength{\parindent}{0pt}
\setlength{\parskip}{0pt}

% Kompiliavimas: pdfLaTeX.
% Spausdinti horizontaliame A4 lape, mastelis 100 % (tikrasis dydis).
% Perkirpus per vidurį gaunami du statūs A5 dydžio lapeliai.

% Tuščias stačiakampis: plotis ir aukštis.
\newcommand{\remelis}[2]{%
  \tikz[baseline={([yshift=-0.6ex]current bounding box.center)}]\draw[line width=0.5pt]
    (0,0) rectangle (#1,#2);}

% Viena užduotis. Numeris ir sąlyga įrašomi kaip du argumentai.
\newcommand{\uzduotis}[2]{%
  \begin{minipage}[t][25mm][t]{\linewidth}
    \textbf{#1.}\enspace #2\par
    \vspace{3mm}
    \remelis{50.4mm}{12mm}\par
  \end{minipage}\par
}

% Vieno A5 lapelio turinys.
% #1 – varianto pavadinimas; #2 – užduočių sąrašas.
\newcommand{\lapelis}[2]{%
  \begin{minipage}[t][190mm][t]{128.5mm}
    \small
    \textbf{Mokinio kodas}\enspace\remelis{32mm}{8mm}%
    \hfill\textbf{#1}\par
    \vspace{6mm}
    {\centering\normalsize\textbf{Bazinių įgūdžių testas 2}\par}
    \vspace{2mm}
    {\centering\small Į rėmelius įrašykite tik atsakymus.\par}
    \vspace{7mm}
    \normalsize
    #2
  \end{minipage}%
}

\newcommand{\stulpeliai}[2]{%
  \begin{minipage}[t]{60mm}#1\end{minipage}\hfill
  \begin{minipage}[t]{60mm}#2\end{minipage}%
}

\newcommand{\papildomaslapas}[4]{%
  \newpage
  \null
  \begin{tikzpicture}[remember picture,overlay]
    \draw[dash pattern=on 2pt off 3pt,line width=0.4pt]
      (current page.north) -- (current page.south);
    \node[anchor=north west,inner sep=0pt]
      at ([xshift=10mm,yshift=-10mm]current page.north west)
      {\lapelis{#1}{#2}};
    \node[anchor=north west,inner sep=0pt]
      at ([xshift=10mm,yshift=-10mm]current page.north)
      {\lapelis{#3}{#4}};
  \end{tikzpicture}%
}

\begin{document}
\null
\begin{tikzpicture}[remember picture,overlay]
  % Punktyrinė kirpimo linija tiksliai A4 lapo viduryje.
  \draw[dash pattern=on 2pt off 3pt,line width=0.4pt]
    (current page.north) -- (current page.south);

  % Kairysis A5 lapelis.
  \node[anchor=north west,inner sep=0pt]
    at ([xshift=10mm,yshift=-10mm]current page.north west) {%
      \lapelis{Kantrybės variantas}{%
        \begin{minipage}[t]{60mm}
          \uzduotis{1}{Išskirkite pilnąjį kvadratą:\\\(2x^2+8x+5\)}
          \uzduotis{2}{Išskirkite pilnąjį kvadratą:\\\(-3x^2+12x-7\)}
          \uzduotis{3}{Atskliauskite:\\\((2x+3)^3\)}
        \end{minipage}\hfill
        \begin{minipage}[t]{60mm}
          \uzduotis{4}{Išskaidykite dauginamaisiais:\\\(27x^3+8\)}
          \uzduotis{5}{Išskaidykite dauginamaisiais:\\\(9-16x^2\)}
          \uzduotis{6}{Išskaidykite dauginamaisiais:\\\(x^2+3x+xy+3y\)}
          \uzduotis{7}{Išskaidykite dauginamaisiais:\\\(4x^2+6xy+6xy+9y^2-8x-12y\)}
        \end{minipage}
      }%
    };

  % Dešinysis A5 lapelis.
  \node[anchor=north west,inner sep=0pt]
    at ([xshift=10mm,yshift=-10mm]current page.north) {%
      \lapelis{Dosnumo variantas}{%
        \begin{minipage}[t]{60mm}
          \uzduotis{1}{Išskirkite pilnąjį kvadratą:\\\(3x^2+12x+16\)}
          \uzduotis{2}{Išskirkite pilnąjį kvadratą:\\\(-2x^2-8x-3\)}
          \uzduotis{3}{Atskliauskite:\\\((3x+2)^3\)}
        \end{minipage}\hfill
        \begin{minipage}[t]{60mm}
          \uzduotis{4}{Išskaidykite dauginamaisiais:\\\(1000x^3-1\)}
          \uzduotis{5}{Išskaidykite dauginamaisiais:\\\(25-4x^2\)}
          \uzduotis{6}{Išskaidykite dauginamaisiais:\\\(x^2+4x+xy+4y\)}
          \uzduotis{7}{Išskaidykite dauginamaisiais:\\\(9x^2+6xy+6xy+4y^2-15x-10y\)}
        \end{minipage}
      }%
    };
\end{tikzpicture}

\papildomaslapas{Drąsos variantas}{%
  \stulpeliai{%
    \uzduotis{1}{Išskirkite pilnąjį kvadratą:\\\(2x^2-12x+17\)}
    \uzduotis{2}{Išskirkite pilnąjį kvadratą:\\\(-4x^2+16x-13\)}
    \uzduotis{3}{Atskliauskite:\\\((4x+1)^3\)}
  }{%
    \uzduotis{4}{Išskaidykite dauginamaisiais:\\\(64x^3+27\)}
    \uzduotis{5}{Išskaidykite dauginamaisiais:\\\(49-9x^2\)}
    \uzduotis{6}{Išskaidykite dauginamaisiais:\\\(x^2+2x+xy+2y\)}
    \uzduotis{7}{Išskaidykite dauginamaisiais:\\\(x^2+4xy+4xy+16y^2-2x-8y\)}
  }%
}{Išminties variantas}{%
  \stulpeliai{%
    \uzduotis{1}{Išskirkite pilnąjį kvadratą:\\\(4x^2+16x+19\)}
    \uzduotis{2}{Išskirkite pilnąjį kvadratą:\\\(-3x^2-6x+2\)}
    \uzduotis{3}{Atskliauskite:\\\((3x+4)^3\)}
  }{%
    \uzduotis{4}{Išskaidykite dauginamaisiais:\\\(125x^3-8\)}
    \uzduotis{5}{Išskaidykite dauginamaisiais:\\\(4-25x^2\)}
    \uzduotis{6}{Išskaidykite dauginamaisiais:\\\(x^2+5x+xy+5y\)}
    \uzduotis{7}{Išskaidykite dauginamaisiais:\\\(16x^2+4xy+4xy+y^2-12x-3y\)}
  }%
}

\papildomaslapas{Teisingumo variantas}{%
  \stulpeliai{%
    \uzduotis{1}{Išskirkite pilnąjį kvadratą:\\\(5x^2-10x+2\)}
    \uzduotis{2}{Išskirkite pilnąjį kvadratą:\\\(-2x^2+12x-17\)}
    \uzduotis{3}{Atskliauskite:\\\((2x+5)^3\)}
  }{%
    \uzduotis{4}{Išskaidykite dauginamaisiais:\\\(8x^3+125\)}
    \uzduotis{5}{Išskaidykite dauginamaisiais:\\\(81-16x^2\)}
    \uzduotis{6}{Išskaidykite dauginamaisiais:\\\(x^2+6x+xy+6y\)}
    \uzduotis{7}{Išskaidykite dauginamaisiais:\\\(4x^2+10xy+10xy+25y^2-2x-5y\)}
  }%
}{Saikingumo variantas}{%
  \stulpeliai{%
    \uzduotis{1}{Išskirkite pilnąjį kvadratą:\\\(2x^2+4x+7\)}
    \uzduotis{2}{Išskirkite pilnąjį kvadratą:\\\(-5x^2-10x-3\)}
    \uzduotis{3}{Atskliauskite:\\\((5x+2)^3\)}
  }{%
    \uzduotis{4}{Išskaidykite dauginamaisiais:\\\(216x^3-1\)}
    \uzduotis{5}{Išskaidykite dauginamaisiais:\\\(36-49x^2\)}
    \uzduotis{6}{Išskaidykite dauginamaisiais:\\\(x^2+7x+xy+7y\)}
    \uzduotis{7}{Išskaidykite dauginamaisiais:\\\(25x^2+5xy+5xy+y^2-10x-2y\)}
  }%
}

\papildomaslapas{Darbštumo variantas}{%
  \stulpeliai{%
    \uzduotis{1}{Išskirkite pilnąjį kvadratą:\\\(3x^2-6x+8\)}
    \uzduotis{2}{Išskirkite pilnąjį kvadratą:\\\(-4x^2-8x+1\)}
    \uzduotis{3}{Atskliauskite:\\\((4x+3)^3\)}
  }{%
    \uzduotis{4}{Išskaidykite dauginamaisiais:\\\(27x^3-64\)}
    \uzduotis{5}{Išskaidykite dauginamaisiais:\\\(1-64x^2\)}
    \uzduotis{6}{Išskaidykite dauginamaisiais:\\\(x^2+8x+xy+8y\)}
    \uzduotis{7}{Išskaidykite dauginamaisiais:\\\(9x^2+12xy+12xy+16y^2-6x-8y\)}
  }%
}{Sąžiningumo variantas}{%
  \stulpeliai{%
    \uzduotis{1}{Išskirkite pilnąjį kvadratą:\\\(4x^2-8x+15\)}
    \uzduotis{2}{Išskirkite pilnąjį kvadratą:\\\(-3x^2+18x-25\)}
    \uzduotis{3}{Atskliauskite:\\\((3x+5)^3\)}
  }{%
    \uzduotis{4}{Išskaidykite dauginamaisiais:\\\(8x^3-27\)}
    \uzduotis{5}{Išskaidykite dauginamaisiais:\\\(100-9x^2\)}
    \uzduotis{6}{Išskaidykite dauginamaisiais:\\\(x^2+9x+xy+9y\)}
    \uzduotis{7}{Išskaidykite dauginamaisiais:\\\(4x^2+2xy+2xy+y^2-12x-6y\)}
  }%
}

\papildomaslapas{Kuklumo variantas}{%
  \stulpeliai{%
    \uzduotis{1}{Išskirkite pilnąjį kvadratą:\\\(2x^2+16x+35\)}
    \uzduotis{2}{Išskirkite pilnąjį kvadratą:\\\(-2x^2+4x+1\)}
    \uzduotis{3}{Atskliauskite:\\\((5x+1)^3\)}
  }{%
    \uzduotis{4}{Išskaidykite dauginamaisiais:\\\(343x^3+8\)}
    \uzduotis{5}{Išskaidykite dauginamaisiais:\\\(121-4x^2\)}
    \uzduotis{6}{Išskaidykite dauginamaisiais:\\\(x^2+10x+xy+10y\)}
    \uzduotis{7}{Išskaidykite dauginamaisiais:\\\(x^2+3xy+3xy+9y^2-4x-12y\)}
  }%
}{Atjautos variantas}{%
  \stulpeliai{%
    \uzduotis{1}{Išskirkite pilnąjį kvadratą:\\\(3x^2+18x+29\)}
    \uzduotis{2}{Išskirkite pilnąjį kvadratą:\\\(-5x^2+20x-17\)}
    \uzduotis{3}{Atskliauskite:\\\((2x+7)^3\)}
  }{%
    \uzduotis{4}{Išskaidykite dauginamaisiais:\\\(729x^3-1\)}
    \uzduotis{5}{Išskaidykite dauginamaisiais:\\\(144-25x^2\)}
    \uzduotis{6}{Išskaidykite dauginamaisiais:\\\(x^2+11x+xy+11y\)}
    \uzduotis{7}{Išskaidykite dauginamaisiais:\\\(16x^2+12xy+12xy+9y^2-20x-15y\)}
  }%
}
\end{document}
